% !TeX root = ../papers/nc-p3.tex
\section{二窓のExt微分と通常の経路空間による反復}
\label{sec:two-window-ext}
以下では前節の$B$自体の矢が$0\to1\to2\to3$となるよう番号を振る。
従って元の二窓からは、順序を逆にし各関係の二文字を反転して入力する。
$X_a,L_a$をこの向きの矢空間と関係空間とし、$S=B_0=k^4$、
$K_{a,p}$を長さ$p$のKoszul交叉とする。$p=0,1,2,3$で
それぞれ$k,X_a,L_a,U$（存在する端点のみ）であり、$U$はこの向きの三次交叉である。
$K_p=\bigoplus_a K_{a,p}$と置く。

\begin{proposition}[四項の計算と新しい辺の表示]\label{prop:two-window-ext-complex}
任意の右$B$加群$M$と頂点$i$について、$\operatorname{Ext}_B^p(DB,M)e_i$は、
次の$p=0,1,2,3$の複体で計算できる。
\[
 D_i^p(M)=\bigoplus_{0\le a\le\min(i,3-p)}
 (e_aBe_i)\otimes K_{a,p}^*\otimes M e_{a+p}.
\]
$T=\operatorname{Ext}_B^3(DB,B)$とすると、消滅を仮定せず任意の$M$について
$\operatorname{Ext}_B^3(DB,M)=M\otimes_B T$である。
$T_B(T)$は$B$に新しい辺$X_3=U^*$を加えた四周期の二次経路代数であり、
新しい辺をまたぐ関係は二つの包含
\[
 U\hookrightarrow L_0\otimes X_2,\qquad U\hookrightarrow X_0\otimes L_1
\]
の係数、すなわち$L_0^*\to X_2\otimes U^*$と
$L_1^*\to U^*\otimes X_0$の像である。
\end{proposition}
\begin{proof}
四層二次代数はKoszulである（補題\ref{lem:four-layer-koszul}）。
二側Koszul分解の項は$B\otimes_S K_p\otimes_S B$で、微分は先頭を左に掛ける項と
末尾を右に掛ける項の$(-1)^p$倍の和である。
この分解の正確性は右端の$B$の次数の降順フィルトレーションでも分かる。
随伴次数では末尾を掛ける微分が消え、一側Koszul複体の直和だけが残る。
増大写像を含めた分解は左$B$加群として分裂するため、左から$DB\otimes_B-$を施すと
右射影分解を得る。$e_i$を左から取って$\operatorname{Hom}_B(-,M)$を施し、
$e_i DB e_a=D(e_aBe_i)$を使うと表示式となる。
微分$D^p\to D^{p+1}$の第一項は$a$を一つ減らし、経路と$K^*$に同じ文字とその双対を
左から掛けて和を取る。第二項は$a$を保ち、$K^*$と$M$の右に対応する文字を掛け、
$(-1)^{p+1}$を付ける。従って写像には矢の積と関係の係数だけが必要である。

$M=B$の複体を全$i$で合わせると両側加群の複体となり、
その項は$B\otimes_S K_p^*\otimes_S B$である（双対は端点を反転する）。
任意の$M$の複体は左から$M\otimes_B-$を施したものなので、最終余核の式が従う。
最終項は端点$3\to0$の$U^*$が生成する自由な相対両側加群である。
長さ2の二項から来る微分ではそれぞれ片側にしか文字を付加できないので、
二つの係数写像の像が全関係を生成する。これをテンソル代数に入れれば主張の表示を得る。
\end{proof}

新辺の出現回数を$m$とすれば、$M_m^{(s)}=e_s T^{\otimes_B m}$の位置$i$成分は
上の四周期経路代数の始点$s$、長さ$4m+i-s$の経路空間そのものである。
負の長さは零とする。これは反復の途中で低次Extが出る場合にも\emph{通常の}テンソル積について成立する。
前の全段が純粋なときに初めて、この複体が次の導来反復の全コホモロジーを計算する。
従って両側射影複体を段ごとに作り直す必要はない。
先頭の左端を固定した経路を一文字ずつ延長し、最後の二次関係の像を割れば$M_m^{(s)}$の
次元と矢作用が得られる。最初の段が全端点で純粋なら、上の両側複体を次数3で終わるように
シフトしたものは$T$の左射影分解でもある。
このとき後続の低次Ext消滅は
$\operatorname{Tor}_{j}^{B}(T^{\otimes_B m},T)=0$ ($1\le j\le3$, $m\ge1$)と同値である。

\paragraph{有限の伝播条件を持つ枝。}
復元した四窓が有限Frobenius条件を満たすとする。
各矢に順序付き基底を選び、同じ長さの語を左辞書順で比較する。
位置$i$の六関係を掃き出して得る六つの先頭二文字の集合を$J_i$とする。
各$i$で、先頭二文字が$J_i$に入り末尾二文字が$J_{i+1}$に入る三文字が
ちょうど四つなら、四窓は全次数でKoszul AS正則である。
実際、次数3の正規語の数は$64-24-24+4=20$であり、
実際の経路商次元も三次交叉から20である。従って正規三文字は独立で、
全ての重なりの差は簡約して零になる。辞書順は固定端点・長さで停止し、左右への連接と両立する。
二次関係間の非自明な重なりは長さ3に尽きるため、ダイヤモンド補題から全正規語が基底となる。
二次PBW代数のKoszul性と有限Frobenius条件がAS性を与える。
逆にこの基底順で二次PBWなら、同じ次数3の次元等式より重なり数は四つである。

この枝では、任意の$s$と$m\ge0$について、各位置に長さ$4m+i-s$の正規語を置き、
矢を付けて簡約するという一つの規則で加群のクラスを記述できる。
$m=0$で全射影を含み、上のテンソル表示で$\operatorname{Ext}^3$により$m$が一つ増える。
PBWから全次数Koszul性が従うので、定理\ref{thm:cut-serre}と帰納法により、
これら全加群の$\operatorname{Ext}^{0,1,2}$は消える。異なる順序や他の枝の被覆は別の問題である。

\paragraph{二回戻る次数8の検算。}
\path{computations/two_window_ext.py}に四項の微分を有理数で実装した。
符号付きOre対照点では四つの始点の一段目が全て純粋で、
二段目の同じ端点でのExt次元は全て$(0,0,0,165)$である。
Boolean境界の始点1,2,3では一段目が純粋であるが、同じ二段目の
$\operatorname{Ext}^2$の次元は順に74,122,122である。
始点0では一段目の純粋性がすでに破れるため、同じ手順の二段目を導来反復の値とはしない。
\path{research/two_window_ext_certificate.json}に端点別の値を置く。
